% ROUND7-REFEREE-POSITIVE-CLOSURE
\section{Round-seven reconstruction: biseam completion, anisotropic response, and an explicit autonomous suspension}
\label{sec:r7-a1}

This section is self-contained and replaces the round-six controlling module.
The point of the reconstruction is that the return section remembers both the
next vertical port and the preceding horizontal port.  Thus no two boundary
germs are collapsed.  The response theory is built on an anisotropic scale for
the invertible baker map rather than on an isotropic Holder space.

\subsection{The biseam graph completion}

Let
\[
0=s_0(a)<s_1(a)<\cdots<s_4(a)=1,
\qquad p_i(a)=s_i(a)-s_{i-1}(a)>0,
\]
and let
\[
0=t_0(a)<t_1(a)<\cdots<t_4(a)=1,
\qquad t_i(a)-t_{i-1}(a)=p_i(a).
\]
Write $I_i=(s_{i-1},s_i)$ and $J_k=(t_{k-1},t_k)$.  The
\emph{biseam completion} is the finite cell complex
\[
 \widehat\Sigma_a=
 \coprod_{1\le i,k\le4}\overline{I_i\times J_k}^{\,g}/\sim,
\]
where two points are identified only when their oriented germs agree in both
coordinate directions.  Consequently an interior point has one cell label,
a point on one seam has two oriented copies, and a crossing of a vertical and
a horizontal seam has four quadrant copies.  We write a point as
$[q,p;i,k]$ whenever its vertical and horizontal germs belong to $I_i$ and
$J_k$.

On a point with vertical label $i$ define
\[
 Q_i(q)=\frac{q-s_{i-1}}{p_i},
 \qquad P_i(p)=t_{i-1}+p_i p.
\]
The next vertical germ $j$ is the unique oriented germ for which
$Q_i(q)\in\overline I_j^{\,g}$; the next horizontal germ is $i$, because
$P_i(p)\in\overline J_i^{\,g}$.  Set
\[
 F_a[q,p;i,k]=[Q_i(q),P_i(p);j,i].
\]

\begin{lemma}[Biseam bijection]
\label{lem:r7-a1-biseam}
The map $F_a:\widehat\Sigma_a\to\widehat\Sigma_a$ is a bijective
stratum-preserving exact symplectic map.  On every closed cell it is affine,
its determinant is one, and its inverse is
\[
 q=s_{i-1}+p_iQ,
 \qquad p=\frac{P-t_{i-1}}{p_i},
\]
where the horizontal germ of the target uniquely supplies the old branch
$i$.  The family is $C^{m+2}$ in $a$ in moving-cell charts whenever the
functions $s_i$ are $C^{m+2}$ and $p_i\ge p_*>0$.
\end{lemma}

\begin{proof}
On the source cell $(i,k)$ the displayed affine map has derivative
$\operatorname{diag}(p_i^{-1},p_i)$ and hence preserves $dq\wedge dp$.
The target horizontal germ equals $i$ even when $P_i(p)$ lies on a horizontal
seam; this is the datum missing from a one-seam completion.  It therefore
selects one inverse branch.  The inverse formula recovers $q,p$, and the
recovered value of $p$ selects the old horizontal germ $k$, including its
oriented side at a seam.  Thus both compositions are the identity on every
quadrant germ.  Since the affine branch primitives differ by exact constants,
the maps paste to an exact symplectic map on the finite cell complex.  In the
coordinates
\[
 u=(q-s_{i-1})/p_i,
 \qquad v=(p-t_{k-1})/p_k,
\]
all incidence maps are fixed and the coefficients are $C^{m+2}$ in $a$.
\end{proof}

\subsection{A complete polyhedral current scale}

Let $\mathcal C_r$ be the set of codimension-$r$ oriented cells of
$\widehat\Sigma_a$, $0\le r\le2$.  For a multi-index $\alpha$ tangential to a
cell and a normal multi-index $\beta$, let
$\operatorname{tr}_{C}^{\alpha,\beta}u$ denote the corresponding current
trace.  The incidence signs are those of the cellular boundary operator, so
corner traces occur as boundaries of face traces and not as independent
formal symbols.

For integers $m\ge0$, $r_u\ge2m+4$, and $r_s\ge2m+4$, define the strong and
weak anisotropic norms by completing cellwise smooth functions for
\[
 \|u\|_{m,+}=
 \sum_{C\in\mathcal C_0}\|u_C\|_{\mathsf B^{r_u,-r_s}}
 +\sum_{\substack{C\in\mathcal C_1\cup\mathcal C_2\\
                   |\alpha|+|\beta|\le m}}
  \omega_C\|\operatorname{tr}_{C}^{\alpha,\beta}u\|_{
  \mathsf B^{r_u-|\alpha|,-r_s-|\beta|}},
\]
\[
 \|u\|_{m,-}=
 \sum_{C\in\mathcal C_0}\|u_C\|_{\mathsf B^{r_u-1,-r_s-1}}
 +\sum_{\substack{C\in\mathcal C_1\cup\mathcal C_2\\
                   |\alpha|+|\beta|\le m}}
  \omega_C\|\operatorname{tr}_{C}^{\alpha,\beta}u\|_{
  \mathsf B^{r_u-|\alpha|-1,-r_s-|\beta|-1}}.
\]
Here $\mathsf B^{r,-s}$ is the completion of Fourier polynomials on a cell
with norm
\[
 \|u\|_{\mathsf B^{r,-s}}
 =\sum_{n,\ell\in\mathbb Z}
  |\widehat u(n,\ell)|(1+|n|)^r(1+|\ell|)^{-s},
\]
using the expanding coordinate in the positive exponent and the stable
coordinate in the negative exponent.  Fixed partition-of-unity changes give
equivalent norms.  The weights $\omega_C$ are chosen recursively so that an
incidence map has norm at most one half of the parent-cell contribution.
Denote the two completions by $\mathscr B_{a,m}^+$ and
$\mathscr B_{a,m}^-$.  The embedding
$\mathscr B_{a,m}^+\hookrightarrow\mathscr B_{a,m}^-$ is compact.

Let $\mathcal L_a$ be the Perron operator of $F_a$ with respect to Liouville
area, acting on all cell and current coordinates by pullback and the cellular
incidence rule.

\begin{theorem}[Anisotropic spectral packet]
\label{thm:r7-a1-anisotropic}
For every finite $m$ there are exponents $r_u,r_s$ and constants
$C<\infty$, $0<\rho<1$, uniform for $a$ in a compact parameter interval, such
that
\[
 \|\mathcal L_a^n u\|_{m,+}
 \le C\rho^n\|u\|_{m,+}+C\|u\|_{m,-}.
\]
The peripheral spectrum on $\mathscr B_{a,m}^+$ consists of the simple
eigenvalue one.  If $\Pi_a u=\nu_a(u)\mathbf1$ and $Q_a=I-\Pi_a$, then
\[
 S_a=(I-\mathcal L_a)^{-1}Q_a
 =\sum_{n\ge0}\mathcal L_a^nQ_a
\]
converges from $\mathscr B_{a,m}^+$ to itself.  Moreover
\[
 a\longmapsto\mathcal L_a
 \quad\hbox{is }C^m\hbox{ as a map }
 \mathscr B_{a,m+j}^+\to\mathscr B_{a,m}^-
\]
for every derivative order $j\le m$ after the moving-cell identifications.
\end{theorem}

\begin{proof}
On branch $i$, unstable frequencies are multiplied by $p_i^{-1}$ under the
map and stable frequencies are multiplied by $p_i$ under the inverse
pullback.  Splitting at frequency $N$ and using $p_i\ge p_*$ gives
\[
 \|\mathcal L_a u\|_{\mathsf B^{r_u,-r_s}}
 \le \rho_0\|u\|_{\mathsf B^{r_u,-r_s}}
 +C_N\|u\|_{\mathsf B^{r_u-1,-r_s-1}},
 \qquad
 \rho_0=\max_i p_i^{\min(r_u,r_s)}<1,
\]
on zero unstable average.  Differentiating cell indicators produces their
oriented face currents; differentiating those produces the signed corner
currents.  Because the cellular boundary squares to zero, all multiple
incidences cancel in compatible pairs.  The recursive choice of $\omega_C$
absorbs the finite triangular incidence matrix, and the same estimate holds
for every mixed trace of total order at most $m$.

Compactness follows by a diagonal Fourier cutoff on each of the finitely many
cells and traces.  Hennion's theorem gives essential spectral radius below
one.  If $\mathcal L_au=\lambda u$, $|\lambda|=1$, the unstable Fourier
estimate makes $u$ constant along unstable leaves; applying the inverse
estimate makes it constant along stable leaves.  The biseam Markov graph is
the full four-symbol graph, hence $u$ is globally constant.  Positivity gives
simplicity.  The resolvent series follows.

In moving-cell coordinates a parameter derivative is a finite sum of an
interior coefficient derivative, a face current, and, after repeated
differentiation, a mixed face/corner current.  Each derivative consumes one
strong derivative and lands in the weak space.  The displayed scale contains
all such currents through order $m$, proving the parameter assertion.
\end{proof}

\begin{theorem}[Finite-order physical response without isotropic-gap input]
\label{thm:r7-a1-response}
Let $G_a\in C^m(\mathscr B_{a,2m+2}^+,\mathbb R)$.  Then
$a\mapsto\nu_a(G_a)$ is $C^m$.  Every derivative is a finite sum of typed
words in
\[
 D^j\mathcal L_a,
 \qquad D^jG_a,
 \qquad S_a,
 \qquad\Pi_a,
\]
where the domain loss is paid by the scale above.  In physical coordinates
the words include all oriented face and corner saltation currents.  The same
quantities are limits of common-root finite differences under the cellwise
quantile coupling.
\end{theorem}

\begin{proof}
Differentiate $\nu_a\mathcal L_a=\nu_a$ and
$\nu_a(1)=1$ on the strong-to-weak scale.  The first derivative is
\[
 D\nu_a(G_a)=\nu_a(DG_a)+
 \nu_a((D\mathcal L_a)S_aQ_aG_a).
\]
The Keller--Liverani resolvent identity on the scale just proved permits
iteration.  Each shape derivative is exactly the cellular current derivative
of the moving partition, so no corner distribution is omitted.  Normal
convergence of the resolvent series and the finite number of incidence types
justify termwise differentiation.  The common-root coupling crosses the same
oriented cells; its distributional derivatives are precisely the saltation
currents, which proves the last assertion.
\end{proof}

\subsection{An explicit autonomous Hamiltonian realization}

For each affine branch choose an exact symplectic isotopy
$R_{i,\tau}$ from the identity to
$(q,p)\mapsto(Q_i(q),P_i(p))$, supported in a Darboux collar assigned to
branch $i$.  The collars and clock intervals are pairwise disjoint.  Let
$K_{a,\tau}$ be the sum of the corresponding compactly supported Hamiltonians,
periodically extended in $\tau\in\mathbb R/\mathbb Z$.  On
\[
 \mathcal M_a=\widehat\Sigma_a\times T^*(\mathbb R/\mathbb Z)
\]
with symplectic form
$\Omega_a=dq\wedge dp+dE\wedge d\tau$, set
\[
 H_a(q,p,\tau,E)=E+K_{a,\tau}(q,p).
\]
The equation $\dot\tau=\partial_EH_a=1$ fixes the return time exactly.
At the end of a clock period the biseam recutting is the identity on physical
coordinates and only changes the cell chart.

\begin{theorem}[Autonomous biseam suspension]
\label{thm:r7-a1-suspension}
The Hamiltonian flow of $H_a$ has first return time one to the section
$\tau=0$ and first return map $F_a$.  It is smooth on every corner stratum,
its impact extensions agree on all oriented face and corner germs, and its
Liouville flux on the section is $dq\,dp$.  The construction is $C^{m+1}$ in
$a$ on every compact parameter interval.
\end{theorem}

\begin{proof}
Hamilton's equations on a branch collar are the prescribed exact isotopy,
while outside the active clock interval they are the identity.  Since the
supports are disjoint, no compatibility equation couples two branch
Hamiltonians.  The biseam incidence rule identifies only equal oriented germs;
the isotopies have matching flat jets at collar boundaries and hence paste on
faces and corners.  The equation $\dot\tau=1$ gives return time one, and the
stroboscopic map is the ordered branch isotopy followed by biseam recutting,
namely $F_a$.  Symplecticity gives the stated flux.  Parameter regularity
follows from the explicit affine generators and a relative Moser correction
with fixed support.
\end{proof}

\subsection{Full-port work and calibration}

Choose a clock interval preceding every branch isotopy and a smooth function
$\rho$ supported there with integral one.  On the incoming vertical germ $i$
put
\[
 \alpha_a=\epsilon_i\rho(\tau)d\tau.
\]
The four incoming components are disjoint in the biseam completion, including
all regular points and all oriented boundary germs.

\begin{proposition}[Exact full-port work]
\label{prop:r7-a1-work}
For every regular return segment entering through port $i$,
\[
 \int_0^1\alpha_a(X_{H_a})dt=\epsilon_i.
\]
Consequently the mechanical work is the additive cocycle
\[
 W_n(x)=\sum_{j=0}^{n-1}\epsilon_{\iota(F_a^jx)}
\]
on the full Liouville section.  No positive-measure core restriction is used.
\end{proposition}

\begin{proof}
During the work interval $K_{a,\tau}=0$, $\dot\tau=1$, and the incoming port
label is constant.  Integration of $\epsilon_i\rho(\tau)d\tau$ gives
$\epsilon_i$.  Concatenation gives the cocycle.
\end{proof}

\begin{theorem}[Conditional mechanical--valuation identification]
\label{thm:r7-a1-calibration}
A smooth cash-additive law-invariant strongly time-consistent certainty
equivalent generated by one utility is either conditional expectation or an
entropic certainty equivalent with a constant coefficient $\vartheta$.
If, in the same work units, its normalized likelihood cocycle is required to
equal the canonical mechanical likelihood generated by $W_n$, then
$\vartheta=\theta$.  Without this compatibility axiom mechanics does not
identify an unrelated preference coefficient.
\end{theorem}

\begin{proof}
Cash additivity yields the translation equation for the utility; differentiating
twice gives constant absolute risk aversion.  The tower property makes its
coefficient time independent.  Equality of the two normalized exponential
likelihoods on the nonconstant full-port cocycle forces
$(\vartheta-\theta)W_n$ to be constant, hence $\vartheta=\theta$.
\end{proof}

\begin{theorem}[Round-seven A1 closure]
\label{thm:r7-a1-main}
The four-port benchmark possesses a genuinely bijective biseam symplectic
return map, an explicit autonomous Hamiltonian suspension with exact return
time, exact work on every regular incoming point, and finite-order response of
arbitrary order on a polyhedral anisotropic current scale.  The construction
uses no isotropic Holder spectral-gap assertion.
\end{theorem}

\begin{proof}
Combine Lemma~\ref{lem:r7-a1-biseam},
Theorems~\ref{thm:r7-a1-anisotropic},
\ref{thm:r7-a1-response}, and
\ref{thm:r7-a1-suspension}, and
Proposition~\ref{prop:r7-a1-work}.  The calibration statement is
Theorem~\ref{thm:r7-a1-calibration}.
\end{proof}
